\section{Methods and Models}
Because the amount of text in input reference documents $(C_i, S_i)$
can be very large
(see Table \ref{wikisum-stats})
it is infeasible to train an end-to-end abstractive model given the memory constraints
of current hardware.
Hence, we first coarsely select a subset of the input using
extractive summarization.
The second stage involves training an abstractive model that generates
the Wikipedia text while conditioning on this extraction. This two-stage
process is inspired by
by how humans might summarize multiple long documents: First highlight
pertinent information, then conditionally generate the summary based on
the highlights.

\subsection{Extractive stage}
We investigate three  extractive methods from the summarization literature,
along with a trivial and cheating method,  to assess the importance of this stage.
For each article, $a_i$ we create a ranked list of paragraphs,
$\{p^i_{R_i(j)}\}$, occurring in $(C_i, S_i)$ where $R_i(j)$ is the rank of the $j$th paragraph $p^i_j$ of $(C_i, S_i)$.
From this we select the first $L$ tokens as input to the second
abstractive stage.


\begin{enumerate}
    \item \textit{Identity}: As a trivial baseline extractor, we simply use the first $L$ tokens of the input.
\item \textit{tf-idf}: A non-trivial ranking is to consider ranking paragraphs as documents in a query-retrieval
problem, where the query is the title of the article, $T(a_i)$. We compute tf-idf \citep{ramos2003using}
for the query, with respect to the documents, $\{p^i_j\}$. That is, we summate
for each word in the query
\[N_w \cdot log (\frac{N_d}{N_{dw}}) \]
where $N_w$, $N_d$, and $N_{dw}$ are the count of the word in the document,
total number of documents, and total number of documents containing the word, 
respectively.

\item \textit{TextRank} \citep{textrank}: A weighted graph is defined
	where text units are nodes and edges are defined
		by a similarity measure based on word overlap.
	An algorithm similar to PageRank \citep{pagerank} is then used to
	compute the ranking of text units. We used paragraphs for the text units.
\item \textit{SumBasic} \citep{sumbasic}: Word frequencies in the input text
	are used to assign scores to words, which are in turn used to score
        sentences. After selecting the best scoring sentence, words in it
        have their scores reduced, and the process is repeated until the desired
	summary length is reached.

\item \textit{Cheating}
    To further demonstrate the quality of extraction on the final performance,
    we implement a cheating extractor that ranks $\{p^i_j\}$ using recall of bigrams
    in the ground truth text:

\begin{equation}
d(p^i_j, a_i) = \frac{bigrams(p^i_j) \cap bigrams(a_i)}{bigrams(a_i)}
\label{cheating-similarity}
\end{equation}
\end{enumerate}


\subsection{Abstractive stage}
\subsubsection{Data representation}
Given the ordered paragraphs $\{p^i_{R_i(j)}\}$,
we derive the raw text input simply as the concatenation of the paragraphs in
order,
the most relevant at the beginning, and prefixed with the title.

We then encode the text using sub-word tokenization similar to \citet{wu2016google}
with a vocabulary size of 32,000 yielding tokenized input, $x_i$:
   \[ text_i = T(a_i) \Vert \{p^i_{R_i(j)}\}  \]
 \[ tokenize(text_i) = x_i = (x_i^1, x_i^2, ..., x_i^{n_i}) \]

For various values of $L$ in experiments, we truncate the tokens to form the input sequence:
  \[ m^L_i = (x_i^1, ... x_i^{min(L, n_i)}) \]

For the output, we use the same vocabulary and tokenization for the Wikipedia
lead text but do not do any truncation across experiments.

Next we describe the abstractive models, $W$, that learn to write articles, $a_i = W(m^L_i)$,
which we treat as a sequence transduction problem from very long input sequences
(up to $L=11000$) to medium output sequences (typically less than 500).

\subsubsection{Baseline Models}

% http://google3/third_party/py/tensor2tensor/models/lstm.py?l=265&rcl=169297690
As a baseline we apply the standard LSTM encoder-decoder with
attention (seq2seq-att) as in \cite{bahdanau2014neural} to this task.
As is typical we train to optimize the maximum-likelihood objective:

\[ y_i = tokenize(a_i) \]
\[\prod_{i=1}^{N}p(y_i | m^L_i) \]

A stronger, more recent baseline that we use is the non-recurrent
Transformer model described in \ref{related_transformer}, which also has symmetric
encoder and decoder modules (T-ED).

\subsubsection{Transformer Decoder (T-D)}
We introduce a simple but effective modification to T-ED for long sequences
that drops the encoder module (almost reducing model parameters by half for a given
hyper-parameter set), combines the input and output
sequences into a single "sentence" and is trained as a standard language model. 

That is, we convert a sequence-transduction example $(m^1, ..., m^n) \mapsto (y^1, ..., y^\eta)$ into
the sentence $(w^1, ..., w^{n+\eta+1}) = (m^1, ..., m^n, \delta, y^1, ..., y^\eta)$, where $\delta$ is a special separator token
 and train a model to predict the next word given the previous ones:
 \[p(w^1, ..., w^{n+\eta}) = \prod_{j=1}^{n+\eta} p(w^i | w^1, ..., w^{j-1}) \]

Since the model is forced to predict the next token in the input, $m$, as well as $y$,
error signals are propagated from both input and output time-steps during training. We also
suspect that for monolingual text-to-text tasks redundant information is re-learned about language in the encoder and decoder. 
 We believe this allows for easier optimization and empirically observe this 
 with longer sequences (see Section \ref{discussion}).
 Note that because of the
 self-attention of the Transformer, when generating the next token, attention
 from both $m$ and $y$ are considered.
 At inference we provide the input sequence, $m_i$, initially, and auto-regressively generate the output, $y_i$, as normal.


\subsubsection{Transformer Decoder with memory-compressed attention (T-DMCA)} \label{t-dmca}
To re-use the terminology used to describe the Transformer,
the attention is a function of a query ($Q$) and set of key ($K$) and value ($V$) pairs.
To handle longer sequences, we modify the multi-head self-attention
of the Transformer to reduce memory usage
by limiting the dot products between $Q$ and $K$ in:
	\[Attention(Q, K, V) = softmax(\frac{QK^T}{\sqrt{d_k}}) V \]

\textbf{Local attention}: Sequence tokens are divided into blocks of similar length and 
    attention is performed in each block independently. As the attention memory cost per block 
    becomes constant, this modification allow us to keep the number of activations linear
    with respect to the sequence length. In our experiments, we choose to have blocks of 256 tokens.

\textbf{Memory-compressed attention}:
After projecting the tokens into the query, key, and value embeddings, we reduce the number of keys
and values by using a strided convolution. The number of queries remains unchanged.
This modification allows us to divide the number of activations by a compression factor.
In our experiments we use convolution kernels of size 3 with stride 3.
In contrast to local attention layers,
which only capture the local information within a block,
the memory-compressed attention layers are able to exchange information
globally on the entire sequence.

These modifications (see Figure \ref{td_model}) allow us in practice to process
sequences 3x in length over the T-D model.  For both local and memory-compressed attention, masking is added to prevent the queries from attending to future keys and values.
Our final architecture is a 5-layer network (LMLML) alternating between local-attention (L) layers and
memory-compressed attention (M) layers (in \citet{vaswani2017attention} it is 6 identical layers).
We also added in some experiments
one mixture of experts (MoE) layer \citep{shazeer2017outrageously}
to increase the network's capacity.


\begin{figure}[h]
\fbox{%
    \includegraphics[width=\textwidth]{AttentionsFigures.png}
}%
	\caption{The architecture of the self-attention layers used in the T-DMCA model.
	Every attention layer takes a sequence of tokens as input and produces
	a sequence of similar length as the output.
	\textbf{Left:} Original self-attention as used in the transformer-decoder.
	\textbf{Middle:} Memory-compressed attention which reduce the number of keys/values.
	\textbf{Right:} Local attention which splits the sequence into individual smaller sub-sequences.
	The sub-sequences are then merged together to get the final output sequence.}
\label{td_model}
\end{figure}
